\section{System Architecture}
\label{sec:arch}
% ====================================================================

\subsection{Overview}

All experiments use the Genesis GPU
simulator~\cite{genesis2024} with the
legged\_gym framework~\cite{rudin2022learning} and
rsl\_rl~\cite{schwarke2025rslrl} on 4$\times$ NVIDIA L40S GPUs.
The robot is a Unitree Go2 quadruped with 12 actuated joints.
Actions $a_t \in \R^{12}$ are target joint positions, updated at 50Hz and converted to torques via a PD controller.
\begin{equation}
  \tau = K_p (a_t - q) - K_d \dot{q},
  \label{eq:pd}
\end{equation} 

The system follows the two-stage privileged distillation pipeline
described in Section~\ref{sec:intro}: a teacher trains with
ground-truth scandot heightmaps~\cite{agarwal2023legged}, then a
student reproduces the teacher's behavior from depth or RGB images
via DAgger (Figure~\ref{fig:pipeline}).

\begin{figure}[t]
  \centering
  % \includegraphics[width=\columnwidth]{figures/pipeline.pdf}
  \fbox{\parbox[c][4.5em][c]{0.9\columnwidth}{\centering
      \textbf{[Pipeline diagram]}\\
      Teacher (scandots) $\to$ Student (depth/RGB)}}
  \caption{Two-stage training pipeline.  The teacher uses privileged
    scandot observations; the student replaces these with either GT depth,
  DA2-inferred depth.}
  \label{fig:pipeline}
\end{figure}

\subsection{Teacher}

The teacher policy $\pi_T$ maps proprioception, privileged physics
parameters, and scandot heightmaps to joint targets.  Three parallel
MLPs encode each observation group into 32-dim latent vectors, which
are concatenated with proprioception and fed through the actor MLP
$\R^{n} \to \R^{12}$ to produce joint position offsets. 

\subsection{Student}
\label{sec:student}

The student replaces the scandot encoder with a visual encoder that
must recover the terrain latent $z \in \R^{32}$ from a depth image (Figure~\ref{fig:student_arch}).
A CNN maps the depth image $d \in \R^{58 \times 87}$ to a 32-dim
feature, which is concatenated with proprioception and compressed by
an MLP.  A GRU integrates these fused features
across timesteps, producing $z$ and a heading estimate
$\hat{y} \in \R^{2}$.  The teacher's actor MLP is reused, taking
$[o_\text{prop};\, z]$ to output 12 joint targets.  During
DAgger distillation, the student minimizes the loss in
Eq.~\ref{eq:dagger_loss} while the teacher's actions are executed
in the environment.

\begin{figure}[t]
  \centering
  \fbox{\parbox[c][5.5em][c]{0.9\columnwidth}{\centering
      \textbf{[Student architecture diagram]}\\
      Depth CNN $\to$ Combo MLP $\to$ GRU $\to$ [latent + yaw]\\
      {}\![proprio; latent] $\to$ Actor MLP $\to$ actions}}
  \caption{Student recurrent depth encoder architecture.  The CNN
    depth backbone and proprioception feed into a GRU that produces
  a terrain latent and heading estimate.}
  \label{fig:student_arch}
\end{figure}

\subsection{Terrain and Curriculum}

The simulator generates parkour terrains with three obstacle
families (stairs, hurdles, gaps) at four difficulty levels.
Gap widths range from 0.24\,m (Row~0) to 0.50\,m (Row~3).  During
teacher training, a curriculum advances robots to harder rows upon
success.
